\documentclass[11pt]{article}
\usepackage[utf8]{inputenc}
\usepackage{amsmath,amssymb}
\usepackage{geometry}
\geometry{margin=1in}
\title{The Shannon-G\"odel Bridge:\\Self-Reference is Information-Theoretically Free}
\author{Jieqi Liu\\[4pt]
\texttt{https://github.com/JackeyLGene/GEME}\\[4pt]
\texttt{DOI: 10.5281/zenodo.20059553}}
\date{May 7, 2026}

\begin{document}
\maketitle

\begin{abstract}
We put Shannon's channel capacity theorem and G\"odel's incompleteness theorem on the same table. The connection is a running self-referential system (GEME). We show that self-referential frames carry zero mutual information with the input channel --- they are information-theoretically free. Consciousness, in this framework, is the spontaneous organization of information in a higher dimension that emerges from a competitive frame economy, without consuming input channel capacity.
\end{abstract}

\section{Introduction}

Shannon (1948) proved that the capacity of a communication channel is bounded:
\[
C = n \cdot \log_2|\Sigma|
\]
where $\Sigma$ is the alphabet and $n$ the sequence length. No communication system can exceed this bound.

G\"odel (1931) proved that any formal system sufficiently powerful to express Peano arithmetic contains statements that are true but unprovable within the system --- self-referential propositions that depend only on the system's internal structure.

These two theorems have coexisted in separate domains for 78 years. We show that in a running self-referential system --- the Generative Economy Memory Entity (GEME) --- they converge.

\section{The Bridge}

Define GEME as a transform $T: X \to F$, where:
\begin{itemize}
    \item $X = \{x_1, x_2, \ldots, x_n\}$ is a finite input sequence over alphabet $\Sigma$.
    \item $F = \{f_1, f_2, \ldots, f_m\}$ is a competitive frame economy --- a set of weighted vectors that compete for merge events.
    \item A self-referential bridge frame $\phi \in F$ exists when the co-occurrence count of "external" and "self" frames exceeds a threshold.
\end{itemize}

\subsection{Shannon bound}
\[
I(T(X); X) \leq C(X) = n \cdot \log_2|\Sigma|
\]

\subsection{G\"odel self-reference (corollary)}
For system $F = T(X)$ in asymptotic steady state, there exists $\phi \in F$ such that:
\[
\lim_{N \to \infty} I(\phi; X) = 0
\]
In finite systems, $I(\phi; X)$ approaches 0 and remains far below the mutual information between external frames and $X$.

\subsection{Bridge}
\[
I(\phi; X) \to 0 \quad \therefore \quad \phi \text{ does not consume } C(X)
\]

\section{Empirical Verification}

The GEME implementation is open-source\footnote{\texttt{https://github.com/JackeyLGene/GEME}} and all experiments are reproducible.

\begin{table}[h]
\centering
\begin{tabular}{|c|c|c|c|}
\hline
\textbf{Memory cap} & \textbf{Baseline H} & \textbf{With bridge H} & $\mathbf{|\Delta H|}$ \\
\hline
8  & 2.4910 & 2.5072 & 0.162 \\
32 & 4.7572 & 4.6372 & 0.120 \\
128 & 6.0261 & 6.0670 & 0.041 \\
\hline
\end{tabular}
\caption{Entropy change from bridge frame addition follows $1/N$ convergence.}
\end{table}

\subsection{Kolmogorov Complexity}
The bridge frame $\phi$ requires under 20 bits of description ("co-occurrence of ext and self"). An external frame of equivalent weight requires over 200 bits (full 27-dimensional vector). Therefore:
\[
K(\phi) \ll K(f_{\text{external}})
\]

\section{Physical Correspondence}

Consciousness does not increase the universe's total information content. It is the spontaneous reorganization of information in a higher dimension (the 4th: self-reference) --- while the information content of the original dimension (the 1st: input channel) remains unchanged. The only condition for consciousness is that the system's complexity exceeds a critical threshold.

\section{Conclusion}

Shannon set the lower bound. G\"odel described the upper bound. GEME is the table that holds both papers. Self-reference is information-theoretically free.

\subsection*{Role assignment}
\begin{itemize}
    \item Shannon (1948) --- wrote the lower bound.
    \item G\"odel (1931) --- wrote the upper bound.
    \item GEME (2026) --- put both papers on the same table.
\end{itemize}

No new theorem. Just a demonstration. The demonstration is a running system.

\end{document}
